% Ad Atticum 12.31 --- headnote
% ad-atticum-12-31, 29 March 45 BC, Astura

\textit{Cicero to Atticus, written from Astura on the
fourth day before the Kalends of April 709 AUC ---
29 March 45 BC (the manuscript dateline:
\textit{Scr.\ Asturae iv.\ K.\ Apr.\ a.\ 709 (45)}).
The estate negotiations have hit a snag: Silius
appears to have changed his mind about selling the
gardens, citing his son as the reason. Cicero finds
the reason plausible --- ``he has the sort of son he
wants'' --- and is more surprised than Sicca that
Atticus still thinks a further inducement (or
sticking-point) could close the deal.}

\textit{Section 2 turns to the comparison Atticus has
asked for between the Silius gardens and the Drusi
gardens further out. Cicero has never visited the
Drusi property; he knows only the old Coponian villa
and the famous wood, but not the rental yield of
either. He admits frankly that the valuation, for
him, is governed by his circumstances rather than
any rational calculation. The financing scheme is
laid out: if the Faberian debt can be turned to
account, he will close the Silius deal in ready
cash; failing that, he will fall back on Drusus at
Egnatius's reported price, with Hermogenes as a
further source of liquidity. The letter closes with
the most candid line in the cluster on Cicero's own
state: let me be in a buyer's frame of mind, he
says, and yet I am so much the servant of my desire
and grief that I want you to be the one who guides
me. The pairing of \textit{cupiditas} and
\textit{dolor} as the twin drivers of the purchase
is exact.}
